%%%%%%%%%%%%%%%%%%%%%%%%%%%%%%%%%%%%%%%
% Trabalho de Sistemas Distribu�dos
%%%%%%%%%%%%%%%%%%%%%%%%%%%%%%%%%%%%%%%

\documentclass[11pt,brazil, a4paper]{article}

%%%%%%%%%%%%%%%%%%%%%%%%%%%%%%%%%%%%%%%
% Preambulo - Configura��o de p�gina
%%%%%%%%%%%%%%%%%%%%%%%%%%%%%%%%%%%%%%%
\usepackage{color}
\usepackage[portuges]{babel}
\usepackage[T1]{fontenc}
\usepackage[ansinew]{inputenc}
\usepackage{array}
\usepackage{colortbl}
\usepackage{epsfig}
\usepackage{multirow}
\usepackage{hyperref}
\usepackage{amsmath}

\usepackage{tikz}
\usetikzlibrary{3d}
\usetikzlibrary{calc,decorations.pathmorphing} % noisy shapes
\usetikzlibrary{fit}					% fitting shapes to coordinates
\usetikzlibrary{backgrounds}	% drawing the background after the foreground
\usetikzlibrary{shadows}
\usetikzlibrary{mindmap}
\usetikzlibrary{positioning,shapes,shadows,arrows}
\usepackage{pgfplots}
\pgfplotsset{compat=1.8}
\usetikzlibrary{arrows,shapes}

\pagestyle{empty}

% Define margem no topo como sendo 1,7 cm
\setlength{\topmargin}{0pt}
\setlength{\voffset}{-1.8cm}

%define margem inferior como 1,5 cm
\setlength{\textheight}{25,6cm}

%define margens laterais como sendo 1,5 cada
\setlength{\oddsidemargin}{0cm}
\setlength{\hoffset}{-1.2cm}
\setlength{\textwidth}{18cm}

%Define campo com sombreamento cinza
\newcolumntype{G}{>{\columncolor[gray]{0.8}[\tabcolsep]\tiny \bf \sf}p{13pt}<{\centering}}
\newcolumntype{T}{>{\tiny \sf}p{13pt}<{\centering}}


\usepackage{xcolor}
\definecolor{verde}{rgb}{0,0.5,0}

\usepackage{listings}
\lstset{
	language=c,
	basicstyle=\ttfamily\normalsize, 
	keywordstyle=\color{blue}, 
	stringstyle=\color{verde}, 
	commentstyle=\color{red}, 
	extendedchars=true, 
	showspaces=false, 
	showstringspaces=false, 
	numbers=left,
	numberstyle=\tiny,
	breaklines=true, 
	breakautoindent=true, 
	captionpos=b,
	xleftmargin=0pt,
}


%%%%%%%%%%%%%%%%%%%%%%%%%%%%%%%%%%%%%%%
% Inicio do documento
%%%%%%%%%%%%%%%%%%%%%%%%%%%%%%%%%%%%%%%
\begin{document}
	
	\selectlanguage{brazil}
	
	\hspace{-0.8cm} \begin{tabular*}{18cm}{p{2cm}p{12cm}@{\extracolsep{\fill}}}
		\multirow{2}{2cm}{\epsfig{file=puclogo_small_bw,width=1.5cm}}
		&  \vspace{1mm} \usefont{T1}{cmss}{bx}{n} {\centering PONTIF�CIA UNIVERSIDADE CAT�LICA DE MINAS GERAIS }\\
		&  \usefont{T1}{cmss}{x}{n} Instituto de Ci�ncias Exatas e Inform�tica \vspace{3mm}  \\
	\end{tabular*}
	
	
	%%%%%%%%%%%%%%%%%%%%%%%%%%%%%%%%%%%%%%%
	% Define caracteristicas da disciplina
	%%%%%%%%%%%%%%%%%%%%%%%%%%%%%%%%%%%%%%%
	\hspace{-0.8cm} \begin{tabular*}{18.5cm}{|p{7cm}@{\extracolsep{\fill}}|p{6cm}|p{1.7cm}|}
		\hline
		\sf Disciplina & \sf Curso  & \sf Per�odo \\
		
		\multicolumn{1}{|r|}{Algoritmos e Estruturas de Dados II} &
		\multicolumn{1}{c|}{Engenharia de Computa��o} &
		\multicolumn{1}{c|}{2$^\circ$} \\ \hline \hline
		\multicolumn{3}{|l|}{\sf Professor} \\
		\multicolumn{3}{|l|}{\hspace{1cm}Kleber Jacques F. de Souza (klebersouza@pucminas.br)} \\ \hline
	\end{tabular*}
	
	
	\vspace{0.1cm}
	
	
	\begin{center}
		{\Large \bf Lista de Exerc�cios - Ordena��o} \\
		%{\normalsize \bf em dupla}
	\end{center}
	

%===================================================================================================


\begin{enumerate}

\item Analise o c�digo abaixo e fa�a:

\begin{enumerate}
	\item execute-o passo a passo no caderno, descrevendo todas as sa�das impressas na tela para o usu�rio. 
	\item Tente descobrir qual o comportamento deste algoritmo de ordena��o.
	\item Este algoritmo se parece com algum algoritmo j� estudado? Se sim, qual?
\end{enumerate}

\begin{lstlisting}
void printVetor(int *vet, int n){
	for(int i=0; i<n; i++)
		printf("%d ",vet[i]);
	printf("\n");
}
void shellSort(int *vet, int n) {
  int h = 1;
  while(h < n) {
    h = 3*h+1;
  }
  while ( h > 1) {
    h /= 3;
    for(int i = h; i < n; i++) {
      int value = vet[i];
      int j = i - h;
      printf("H=%d, Compara se %d < %d\n",h,value,vet[j]);
      while (j >= 0 && value < vet[j]) {
        printf("H=%d, Troucou: %d %d\n",h,vet [j + h],vet[j]);
        vet [j + h] = vet[j];
        j -= h;
      }
      vet [j + h] = value;
      printVetor(vet,n);
    }
  }
}
int main(int argc, char *argv[]){
	int v[] = {15,14,13,12,11,10,9,8,7,6,5,4,3,2,1};
	int n = sizeof(v)/sizeof(int);	
	printf("Antes de Ordenar:\n");
	printVetor(v,n);
	shellSort(v,n);
	printf("Depois de Ordenar:\n");
	printVetor(v,n);	
	return 0;
}

\end{lstlisting}
 

\end{enumerate}






%===================================================================================================

\end{document}